\section{Experiments}
\label{experiments}

\subsection{Research Questions}
The experiments evaluate whether the proposed workflow improves requirement
conflict detection and consistency verification.
\begin{itemize}[leftmargin=*]
  \item \textbf{RQ1:} How accurately does the workflow detect conflicts in natural
  language requirements?
  \item \textbf{RQ2:} How much does knowledge-guided analysis improve consistency
  verification?
  \item \textbf{RQ3:} How does multi-agent collaboration affect false positives and
  explanation quality?
  \item \textbf{RQ4:} To what extent can human-in-the-loop refinement reduce manual
  verification effort?
\end{itemize}

\subsection{Experimental Setup}
The evaluation will compare the proposed workflow with simpler requirement checking
settings under the same input requirements and annotation criteria. The setup should
specify agent prompts, knowledge sources, conflict taxonomy, annotation protocol,
and human verification procedure.

\subsection{Datasets}
The study will use requirement sets containing manually annotated conflicts and
non-conflicts. Dataset reporting should include the application domain, number of
requirements, conflict categories, annotation process, and inter-annotator agreement
where applicable.

\subsection{Baselines}
Baselines should be selected to isolate the effect of collaboration and knowledge
guidance:
\begin{itemize}[leftmargin=*]
  \item A single-agent requirement conflict checker.
  \item A requirement checker without external requirement knowledge.
  \item A rule-only consistency checker where applicable.
\end{itemize}

\subsection{Evaluation Metrics}

\subsubsection{Conflict Detection Accuracy}
Conflict detection accuracy measures whether the workflow correctly distinguishes
conflicting and non-conflicting requirement cases.

\subsubsection{Consistency Verification Accuracy}
Consistency verification accuracy measures whether detected conflicts are assigned
to the correct consistency issue and conflict category.

\subsubsection{False Positive Rate}
False positive rate measures how often the workflow reports conflicts that human
annotators consider acceptable or non-conflicting.

\subsubsection{Human Verification Reduction}
Human verification reduction measures the decrease in manual review effort after
agent-based filtering, explanation, and refinement.

\subsection{Main Results}
This subsection will report the main quantitative results once the annotated
requirement data and baselines are finalized. No synthetic performance numbers are
introduced in the paper structure.

\subsection{Ablation Study}
The ablation study should evaluate the contribution of major components, such as
knowledge guidance, constraint rules, role-separated agents, and human feedback.

\subsection{Case Study}
The case study should present a small requirement set with a representative conflict,
the workflow output, the explanation, and the human refinement decision.
